
\chapter*{Taller autónomo con reporte técnico breve}
\phantomsection\addcontentsline{toc}{chapter}{Taller autónomo y reporte técnico}\label{mod:taller}
\begin{taller}[Finalidad y entrega]
Demostrar, con cálculos y archivos regenerables, cómo las decisiones de
representación afectan a un pequeño modelo de lenguaje. Entregar un
\textbf{reporte de 2--3 páginas} (objetivo, datos, métodos, resultados,
incertidumbre, limitaciones y relación con IC10), más scripts y CSV.
\textbf{Todo es trabajo autónomo}; no se presuponen las horas del laboratorio
dentro de la clase presencial. En S01 se conservan como hito de investigación
\textbf{dos áreas de interés}, cada una con un párrafo.
\end{taller}
\section*{Actividad 1 [Representación]: una frase, cuatro objetos}
Segmente ``el gato mira al perro'' en palabras y caracteres; construya
unigramas, bigramas y trigramas; contraste bolsa de palabras con secuencia y
escriba el vector one-hot del token \texttt{gato}. Repita el primer paso con
una frase técnica. \emph{Autoverificación}: cinco palabras, cuatro pares,
tres trigramas en la primera oración.
\section*{Actividad 2 [Cálculo]: del conteo a la predicción}
Con \texttt{abacabbacb} construya la matriz de transiciones, derive el
estimador MLE e introduzca $\alpha=1$. Para \texttt{abca}, calcule la
probabilidad de transición y la conjunta con $p(X_1=a)=1/3$.
\emph{Autoverificación}: $P(b\mid a)=3/7$; probabilidad conjunta $1/105$.
\section*{Actividad 3 [Demostración]: medir lo que se predice}
Demuestre la identidad entre entropía cruzada y KL con las condiciones de
soporte; calcule KL de $(1/2,1/4,1/4)$ respecto de
$(1/4,1/2,1/4)$. Para las probabilidades observadas
$(1/2,1/4,1/8,1/2)$, calcule PPL. Compare dos tokenizadores con
$(T_A,\mathrm{PPL}_A)=(100,3.2)$ y $(T_B,\mathrm{PPL}_B)=(30,25)$
sobre un mismo texto de 100 caracteres, indicando qué supuestos permiten
interpretar el BPC. \emph{Autoverificación}: KL $=0.25$ bits y
PPL $=2^{7/4}$.
\section*{Actividad 4 [Algoritmos]: BPE, R y Python}
Ejecute seis fusiones BPE sobre la tabla \emph{casa, casas, cosa, cosas,
caso} con las frecuencias y reglas indicadas en la sección~1.1.2 del libro. Ejecute después:
\begin{verbatim}
python codigo/s01.py
python codigo/experimento_s01.py
python codigo/ic10_tokenizadores.py
Rscript codigo/s01.R
Rscript codigo/experimento_s01.R
Rscript codigo/ic10_tokenizadores.R
\end{verbatim}
Compruebe las sumas de las filas, las longitudes BPE y las seis filas
pareadas del CSV. Compare también, sobre las mismas seis frases
de prueba, el BPE general con el tecnico mediante 20 fusiones en ambos;
produzca \texttt{resultados\_ic10\_tokenizacion.csv} y explique las diferencias.
Si no tiene R, indique \textbf{PENDIENTE} en el reporte,
sin inventar que los números se reprodujeron en ambos lenguajes.
\section*{Actividad 5 [Aplicación e investigación]: entregar evidencia}
Use \nolinkurl{experimentos/corpus_didactico.csv} para reconstruir
\nolinkurl{experimentos/resultados_python.csv}; complete la
\textbf{plantilla de reporte} en \texttt{reportes/reporte\_s01.tex}.
Incluya la tabla BPE con tokens por palabra y distinga la comparación
de tokenizadores de la comparación de predictores por caracteres.
Interprete la diferencia media de BPC y el intervalo bootstrap por
\emph{documento}, explicando por qué ni el tamaño de muestra ni el origen
del corpus permiten inferencia para textos industriales reales.
Proponga una versión de IC10 que sí sea falsable con documentos y
comparaciones externas. Finalmente, formule dos áreas de interés,
con un párrafo justificado por cada una.
\begin{verificacion}[Lista de control antes de entregar]
\begin{enumerate}
\item Los datos no se mezclan entre ajuste y prueba; se identifica cada
unidad de observación.
\item Se reportan todas las probabilidades de la cadena, incluido EOS,
y se cuenta explícitamente la unidad del denominador BPC.
\item Cada resultado se etiqueta como \emph{calculado}, \emph{ejecutado
con datos sintéticos} o \emph{pendiente de validación empírica}.
\item Se entregan código, CSV, versión de lenguaje, interpretación,
limitaciones y dos áreas de interés.
\end{enumerate}
\end{verificacion}
\section*{Extensiones opcionales, fuera del núcleo S01}
Calcule intervalo de confianza mediante bootstrap jerárquico con varios
documentos por fuente; incorpore una evaluación BPC de modelos probabilisticos sobre tokens
BPE con distribuciones comparables; compare
vocabularios fijados por entrenamiento con distintos tamaños; pruebe
suavizado $\alpha\in\{1/2,1,2\}$ exclusivamente en validación; diseñe una
réplica externa para dos dominios técnicos con control de licencias.
